\thispagestyle{empty}
\vspace*{30mm}
\begin{center}
{\fontsize{24}{28}\selectfont\itshape La Primera Regla\\[1mm] del Archivo}\\[8mm]
{\color{oxido}\rule{18mm}{0.4pt}}\\[8mm]
{\scshape Una novela corta de Ceres}\\[30mm]
{\small\itshape Prólogo, diez capítulos, epílogo\\ y unos apéndices para quien quiera\\ comprobar las cuentas}
\vfill
{\footnotesize\color{gris} Ficción didáctica basada en el paper\\
\emph{A Fast Wait-Free Solution to Read-Reclaim Races in Reference Counting}\\
(Euro-Par 2024, LNCS 14803, pp.~103--118)\\ y en su explicación detallada en español.}
\end{center}
\clearpage

\thispagestyle{empty}
\vspace*{40mm}
\begin{center}
{\scshape\color{oxido}\addfontfeature{LetterSpace=12} Reglamento General del Archivo de Ceres}\\[2mm]
{\footnotesize\itshape Artículos fundacionales, grabados sobre la puerta de la Bóveda}\\[10mm]
\begin{minipage}{0.8\textwidth}
\setlength{\parskip}{4mm}
\noindent\textsc{Primera Regla.} Ningún cristal será fundido mientras alguien lo esté leyendo.

\noindent\textsc{Segunda Regla.} Ningún lector esperará sin término, salvo que ello contravenga la Primera Regla.

\noindent\textsc{Tercera Regla.} El Archivo no guardará lo que ya nadie lee, salvo que ello contravenga la Primera o la Segunda Regla.
\end{minipage}
\end{center}
\vfill
\begin{flushright}
\footnotesize\itshape «Las tres reglas son fáciles de cumplir por separado.\\ Lo difícil, como con casi todo, es cumplirlas a la vez.»\\[1mm]
\upshape --- Dra.\ Inés Calderón, \emph{Lecciones en la Bóveda}
\end{flushright}
\clearpage

\thispagestyle{empty}
\section*{Personajes}
\begin{description}[leftmargin=0pt,itemsep=3mm,font=\scshape\color{oxido}]
\item[Inés Calderón] Matemática de sistemas del Instituto de Ceres. Nacida en Lima, en la Tierra. Tiene la costumbre de resolver los problemas en la pizarra antes de que nadie termine de plantearlos.
\item[Siete] Robot lector del Archivo, modelo LECTOR, número de serie siete. Literal, preciso y bastante más curioso de lo que su fabricante previó.
\item[Aurelio Mendieta] Comisionado de la Autoridad de Navegación del Cinturón. Quiere un culpable, y lo quiere antes de la cena.
\item[Benito Arce] Archivero de turno. Todos le llaman Beto. Cambió un cartel la noche del incidente, y eso lo ha convertido en el principal sospechoso.
\item[Helga Sørensen] Directora del Archivo y Gran Maestre del Gremio de la Llave. Cree en los candados como otros creen en los santos.
\item[Emeric Loew] Historiador del Archivo, muy viejo, muy sabio y muy dormido, salvo cuando alguien comete un error histórico.
\item[ANA-3] Robot navegante del carguero \emph{Perséfone}. Obedeció todas las reglas y aun así estuvo a punto de estrellar la nave.
\end{description}
\clearpage
